\documentclass[12pt]{scrartcl}
\usepackage[hagavi]{azzam}

% =====================================================================
% Sumber soal (hanya soal yang disalin, solusi tidak ikut disalin):
%   - trig-bash-sunday-june-21.pdf (Sunday Online Class, Trig Bash, June 21, 2020)
%   - bashing-sun-june-28.pdf (Sunday Online Class, Bashing, June 28, 2020)
% Soal yang ada solusinya di PDF asli -> section Example.
% Soal tanpa solusi (homework) -> section Practice.
% Hint Example diringkas dari solusi di PDF asli; hint Practice P1--P7
% disalin dari halaman hint PDF June 21. Soal P8--P16 tidak punya hint
% di PDF asli.
% =====================================================================

\title{Trigonometric Bashing}

\begin{document}
\maketitle

\renewcommand*\contentsname{Daftar Isi}
\tableofcontents

\newpage

\section{Example}

\begin{enumerate}[label=E\arabic*.,leftmargin=*,itemsep=0.6em]

% Sumber: trig-bash-sunday-june-21.pdf No. 1 (dengan solusi)
\item Let $ABCD$ be a convex quadrilateral. Let diagonals $AC$ and $BD$ intersect at $P$. Let $PE$, $PF$, $PG$ and $PH$ be the altitudes from $P$ to the sides $AB$, $BC$, $CD$ and $DA$ respectively. Show that $ABCD$ has an incircle if and only if
\[
\frac{1}{PE} + \frac{1}{PG} = \frac{1}{PF} + \frac{1}{PH}.
\]

% Sumber: trig-bash-sunday-june-21.pdf No. 2 (dengan solusi)
\item Consider an acute-angled triangle $ABC$. Let $P$ be the foot of the altitude of triangle $ABC$ issuing from the vertex $A$, and let $O$ be the circumcentre of triangle $ABC$. Assume that $\angle C \ge \angle B + 30^\circ$. Prove that $\angle A + \angle COP < 90^\circ$.

% Sumber: trig-bash-sunday-june-21.pdf No. 3 (dengan solusi)
\item Let $ABC$ be an acute triangle with $O$ as its circumcentre. Line $AO$ intersects $BC$ at $D$. Points $E$ and $F$ are on $AB$ and $AC$ respectively such that $A$, $E$, $D$, $F$ are concyclic. Prove that the length of the projection of line segment $EF$ on side $BC$ does not depend on the positions of $E$ and $F$.

% Sumber: trig-bash-sunday-june-21.pdf No. 4 (dengan solusi)
\item Let $A_1$ be the center of the square inscribed in acute triangle $ABC$ with two vertices of the square on side $BC$. Thus one of the two remaining vertices of the square is on side $AB$ and the other is on $AC$. Points $B_1$, $C_1$ are defined in a similar way for inscribed squares with two vertices on sides $AC$ and $AB$, respectively. Prove that lines $AA_1$, $BB_1$, $CC_1$ are concurrent.

% Sumber: trig-bash-sunday-june-21.pdf No. 5 (dengan solusi)
\item Suppose $ABCD$ is a cyclic quadrilateral and $r_1$, $r_2$, $r_3$, $r_4$ are respectively the inradii of triangles $ABC$, $ADC$, $BAD$, $BCD$. Show that $r_1 + r_2 = r_3 + r_4$.

% Sumber: bashing-sun-june-28.pdf No. 1 (dengan solusi)
\item Let $ABC$ be a triangle with $\angle BAC = 60^\circ$. Let $AP$ bisect $\angle BAC$ and let $BQ$ bisect $\angle ABC$, with $P$ on $BC$ and $Q$ on $AC$. If $AB + BP = AQ + QB$, what are the angles of the triangle?

% Sumber: bashing-sun-june-28.pdf No. 2 (dengan solusi)
\item Let $ABC$ be a right-angle triangle with $\angle B = 90^\circ$. Let $D$ be a point on $AC$ such that the inradii of the triangles $ABD$ and $CBD$ are equal. If this common value is $r'$ and if $r$ is the inradius of triangle $ABC$, prove that $1/r' = 1/r + 1/BD$.

% Sumber: bashing-sun-june-28.pdf No. 3 (dengan solusi)
\item In an acute triangle $ABC$, let $O$, $G$, $H$ be its circumcentre, centroid and orthocentre, respectively. Let $D \in BC$, $E \in CA$ such that $OD \perp BC$, $HE \perp CA$. Let $F$ be the midpoint of $AB$. If the triangles $ODC$, $HEA$, $GFB$ have the same area, find all the possible values of $\angle C$.

% Sumber: bashing-sun-june-28.pdf No. 4 (dengan solusi)
\item Suppose that $a > b > c > d$ are positive integers such that
\[
ac + bd = (b + d + a - c)(b + d - a + c).
\]
Prove that $ab + cd$ cannot be a prime.

\end{enumerate}

\section{Practice}


\begin{enumerate}[label=P\arabic*.,leftmargin=*,itemsep=0.6em]

% Sumber: trig-bash-sunday-june-21.pdf homework No. 1 (tanpa solusi, ada hint)
\item $ABCDE$ is a convex pentagon satisfying $\angle ABC = 120^\circ$, $\angle CDE = 60^\circ$, $AB = BC$, $CD = DE$, and $BD = 2$. Find the area of $ABCDE$.

% Sumber: trig-bash-sunday-june-21.pdf homework No. 2 (tanpa solusi, ada hint)
\item Let $ABC$ be a triangle with $\angle A = 90^\circ$ and $AB < AC$. Let $AD$ be the altitude from $A$ onto $BC$. Let $P$, $Q$ and $I$ denote respectively the incentres of triangles $ABD$, $ACD$ and $ABC$. Prove that $AI \perp PQ$ and $AI = PQ$.

% Sumber: trig-bash-sunday-june-21.pdf homework No. 3 (tanpa solusi, ada hint)
\item Let $I$ be the incentre of triangle $ABC$ and let ray $AI$ meet the circumcircle of $ABC$ at $D$. Denote the feet of the perpendiculars from $I$ to lines $BD$ and $CD$ by $E$ and $F$, respectively. If $IE + IF = \tfrac12 AD$, calculate $\angle BAC$.

% Sumber: trig-bash-sunday-june-21.pdf homework No. 4 (tanpa solusi, ada hint)
\item Let $ABC$ be an acute triangle with altitudes $AD$, $BE$ and $CF$, and let $O$ be the center of its circumcircle. Show that the segments $OA$, $OF$, $OB$, $OD$, $OC$, $OE$ dissect the triangle $ABC$ into three pairs of triangles that have equal areas.

% Sumber: trig-bash-sunday-june-21.pdf homework No. 5 (tanpa solusi, ada hint)
\item Suppose $AB$ and $CD$ are two non-intersecting chords of a circle. Take any point $P$ on the arc $AB$ not containing $C$, $D$. Let $PC$ and $PD$ intersect $AB$ at $Q$, $R$ respectively. Show that $\dfrac{AQ \cdot RB}{QR}$ is constant (irrespective of position of $P$).

% Sumber: trig-bash-sunday-june-21.pdf homework No. 6 (tanpa solusi, ada hint)
\item Let $ABCD$ be a cyclic quadrilateral inscribed in a circle $\Gamma$. Let $E$, $F$, $G$, $H$ be the midpoints of arcs $AB$, $BC$, $CD$, $AD$ of $\Gamma$, respectively. Suppose that $AC \cdot BD = EG \cdot FH$. Show that $AC$, $BD$, $EG$, $FH$ are all concurrent.

% Sumber: trig-bash-sunday-june-21.pdf homework No. 7 (tanpa solusi, ada hint)
\item Let $H$ be the orthocentre of an acute-angled triangle $ABC$. The circle $\Gamma_A$ centred at the midpoint of $BC$ and passing through $H$ intersects the sideline $BC$ at points $A_1$ and $A_2$. Similarly, define the points $B_1$, $B_2$, $C_1$ and $C_2$. Prove that the six points $A_1$, $A_2$, $B_1$, $B_2$, $C_1$ and $C_2$ are concyclic.

% Sumber: bashing-sun-june-28.pdf homework No. 1 (tanpa solusi, tanpa hint)
\item Let $ABC$ be a triangle with $\angle A = 90^\circ$. Points $D$ and $E$ lie on sides $AC$ and $AB$, respectively, such that $\angle ABD = \angle DBC$ and $\angle ACE = \angle ECB$. Segments $BD$ and $CE$ meet at $I$. Determine whether or not it is possible for segments $AB$, $AC$, $BI$, $ID$, $CI$, $IE$ to all have integer lengths.

% Sumber: bashing-sun-june-28.pdf homework No. 2 (tanpa solusi, tanpa hint)
\item In $\triangle ABC$, all the angles are greater than $30^\circ$. Suppose that $P$ is the centre of a circle that cuts the sides $BC$, $CA$, $AB$ at $T$, $Q$, $K$, $L$, $M$, $N$ respectively. If the triangles $PQK$, $PML$ and $PNT$ are all equilateral, show that
\begin{enumerate}[(a)]
\item the radius of the circle is given by $\dfrac{2abc}{a^2+b^2+c^2+4S\sqrt{3}}$ where $S$ is the area of $\triangle ABC$ and $a$, $b$, $c$ are the side lengths (as usual),
\item $a \cdot PA = b \cdot PB = c \cdot PC$.
\end{enumerate}

% Sumber: bashing-sun-june-28.pdf homework No. 3 (tanpa solusi, tanpa hint)
\item Let $ABCD$ be a convex quadrilateral in which $\angle BAC = 50^\circ$, $\angle CAD = 60^\circ$ and $\angle BDC = 25^\circ$. If $E$ is the point of intersection of $AC$ and $BD$, find $\angle AEB$.

% Sumber: bashing-sun-june-28.pdf homework No. 4 (tanpa solusi, tanpa hint)
\item For $a, b, c > 0$, prove the inequality $\sqrt{a^2-ab+b^2} + \sqrt{b^2-bc+c^2} \ge \sqrt{a^2+ac+c^2}$. Furthermore, show that equality holds if and only if $1/b = 1/a + 1/c$.

% Sumber: bashing-sun-june-28.pdf homework No. 5 (tanpa solusi, tanpa hint)
\item In a cyclic quadrilateral $ABCD$, $AB = a$, $BC = b$, $CD = c$, $\angle ABC = 120^\circ$ and $\angle ABD = 30^\circ$. Prove that, (i) $c \ge a + b$, and (ii) $|\sqrt{c+a} - \sqrt{c+b}| = \sqrt{c-a-b}$.

% Sumber: bashing-sun-june-28.pdf homework No. 6 (tanpa solusi, tanpa hint)
\item Let $A_1A_2\ldots A_n$ be a convex polygon, $n \ge 4$. Prove that $A_1A_2\ldots A_n$ is cyclic if and only if to each vertex $A_j$ one can assign a pair $(b_j, c_j)$ of real numbers, $j = 1, 2, \ldots, n$, so that $A_iA_j = b_jc_i - b_ic_j$ for all $i$, $j$ with $1 \le i < j \le n$.

% Sumber: bashing-sun-june-28.pdf homework No. 7 (tanpa solusi, tanpa hint)
\item Let $ABC$ be a triangle with $AB = AC$. The angle bisectors of $\angle CAB$ and $\angle ABC$ meet the sides $BC$ and $CA$ at $D$ and $E$, respectively. Let $K$ be the incentre of triangle $ADC$. Suppose that $\angle BEK = 45^\circ$. Find all possible values of $\angle CAB$.

% Sumber: bashing-sun-june-28.pdf homework No. 8 (tanpa solusi, tanpa hint)
\item Let $ABC$ be an acute angled triangle. Let $D$, $E$, $F$ be points on $BC$, $CA$, $AB$ such that $AD$ is the median, $BE$ is the internal bisector and $CF$ is the altitude. Suppose that $\angle FDE = \angle C$, $\angle DEF = \angle A$ and $\angle EFD = \angle B$. Show that $ABC$ is equilateral.

% Sumber: bashing-sun-june-28.pdf homework No. 9 (tanpa solusi, tanpa hint)
\item Let $D$, $E$, $F$ be points on the sides $BC$, $CA$, $AB$ respectively of a triangle $ABC$ such that $BD = CE = AF$ and $\angle BDF = \angle CED = \angle AFE$. Show that $\triangle ABC$ is equilateral.

\end{enumerate}

\newpage
\section{Hints}

Putar halaman ini 180 derajat untuk membaca hintnya.

\begin{center}
\rotatebox{180}{%
\begin{minipage}{0.96\linewidth}
\small
\begin{enumerate}[label=E\arabic*.,leftmargin=*,itemsep=0.25em]
\item Relate the altitudes to the sides by areas ($PE = pq\sin\theta/a$, etc.), use Pitot ($a+c = b+d$), and compare $(a+c)^2 = (b+d)^2$ with the altitude condition via the cosine rule.
\item Reduce to $PC < OP$, i.e. $BC \cdot PC < R^2$ or $4\sin A\sin B\cos C < 1$, then use $2\sin B\cos C = \sin A - \sin(C-B)$.
\item Write $XY = BC - (BX + YC)$ with $BX = EB\cos B$, $YC = CF\cos C$; power of a point for $BE$, $CF$; sine rule in $ADB$ and $ADC$.
\item Let $AA_1 \cap BC = X$ and the two parts of $PQ$ be $u$, $v$; get $BX/XC$ from $PQ \parallel BC$ (or from area ratios with $\sin(B+45^\circ)$, $\sin(C+45^\circ)$), then Ceva.
\item Use $r = (s-a)\tan\frac{A}{2}$ for $ABC$ and $ADC$, write the sides with $R\sin$, and simplify to the symmetric sum $R(\cos\alpha+\cos\beta+\cos\gamma+\cos\delta-2)$.
\item Set $\angle ABQ = x$; sine rule in $ABP$ and $ABQ$ gives one equation in $x$; sum to product, then product to sum, then factor with the $\cos\frac{A+B}{2}$ average trick.
\item Similarity gives $r'/r = (AP+QC)/AC$; write $AP = s_1-BD$, $QC = s_2-BD$; inradius $=$ area $/$ semiperimeter; $r = s-b$ for a right triangle.
\item Compute $[ODC]$, $[HEA]$, $[GFB]$ in $R$, $A$, $B$, $C$; equating gives $\tan A = 4\sin C\cos C$ and then an equation only in $C$ with $t = \sin^2 C$.
\item Rewrite as $a^2+c^2-ac = b^2+d^2+bd$ and draw it with $60^\circ$/$120^\circ$ cosine rules as a cyclic quadrilateral; cosine rule for $BD$ plus Ptolemy; order $ab+cd > ac+bd > ad+bc$ from $a > b > c > d$.
\end{enumerate}
\begin{enumerate}[label=P\arabic*.,leftmargin=*,itemsep=0.25em]
\item Sine rule, cosine rule, and the formula $\tfrac12 bc\sin A$ for area.
\item Some angle chase and the formula $AH = 2R\cos A$.
\item Sine rule and Ptolemy's theorem.
\item Formula $\tfrac12 bc\sin A$ for area and other results, such as $CD = b\cos C$.
\item Sine rule and clever elimination.
\item Show that $EF = 2R\sin(D/2)$. Apply Ptolemy's theorem. Finally, use equality case of some inequality.
\item First guess what will be the centre of the circle. Apply Apollonius' theorem. Find parallelograms in the figure.
\end{enumerate}
\end{minipage}%
}
\end{center}

% \section{References}
% \url{https://drive.google.com/file/d/1rM1pop8w74PcS8JyK-PITFJWswgvFZTY/view?usp=drive_link}\\
% \url{https://drive.google.com/file/d/1g-gigP0dWzQAt8dP6sPR5OwWUdaz43OU/view?usp=drive_link}

\end{document}
